% =============================================================================
% 05-process-concept-pcb-notes.tex — Lecture Notes for Lecture 5:
% Process Concept and the Process Control Block
%
% DERIVED ARTIFACT. Source of truth: Slides/CS502/05-process-concept-pcb.tex.
% Generated by /lecture-notes; do not hand-edit content independently of the
% Beamer source without re-running /qa-notes afterward (single-source-of-truth.md).
% Structured as a textbook chapter (numbered sections/figures/examples,
% topic-organized rather than a frame-by-frame transcript).
%
% Anchor reading: Silberschatz et al., OSC 8e, Ch. 3 §3.1 (pp. 101-105);
%   Stallings, OS 6e, Ch. 3 §3.1 (p.108), §3.2 (p.111), §3.3 (p.126);
%   Bovet & Cesati, ULK 3e, Ch. 3 (pp. 79-130). Page-verified via
%   master_supporting_docs/CS502/supporting_books/*/index.md.
%
% Compile with: cd Notes/CS502 &&
%   TEXINPUTS=../../Preambles: xelatex -interaction=nonstopmode 05-process-concept-pcb-notes.tex  (x3)
% =============================================================================

\documentclass[11pt]{article}
\usepackage[margin=1in]{geometry}
\usepackage{array}
\input{../../Preambles/header}
\coursecode{CS502}

\renewcommand{\thesection}{5.\arabic{section}}
\renewcommand{\thefigure}{5.\arabic{figure}}
\renewcommand{\thetable}{5.\arabic{table}}

\newtheorem{example}{Example}
\renewcommand{\theexample}{5.\arabic{example}}

\newenvironment{definitionbox}[1]{%
  \par\noindent\textbf{Definition (#1).}\ \itshape
}{\par\vspace{0.3em}}

\title{Process Concept and the Process Control Block\\\large Lecture Notes --- Lecture 5\\
  \normalsize What Is a Running Program?}
\author{Auqib Hamid Lone\\[0.3em]
  \emph{Department of Computer Science \& Engineering}, \emph{GCET Ganderbal}}
\date{\today}

\begin{document}
\maketitle

\section{One File, Two Runs}

Compile \texttt{counter.c} into a single executable file, \texttt{./counter},
and run it in two terminals at the same time. Each copy counts upward from 1 at
its own pace. There is \emph{one} file on disk, yet the operating system is
clearly running \emph{two} things. The question that organizes this chapter is:
what exactly are those two things, and what must the kernel record about each?

Three observations frame the answer. First, the file on disk does not change
while it runs; it is just bytes. Second, each run has its own counter value, its
own position in the code, and its own stack. Third, the OS must remember all of
this for each run separately. We will call the two runs \texttt{P1} and
\texttt{P2} and follow them through the chapter, asking at every step what the
kernel must record about each.

The chapter answers four questions in turn: \key{what is a process} (a program
in execution, not the program file); \key{what does a process look like in
memory} (text, data, heap, and stack); \key{what states can a process be in}
(new, ready, running, waiting, terminated); and \key{how does the kernel
remember a process} (the process control block, or PCB). The previous lecture
organized the kernel itself; this one names the unit of work the kernel
manages, and the next lecture watches that unit move between queues.

\section{Program versus Process}

\noindent\muted{\textit{Reading: Silberschatz et al., Ch.~3 \S3.1, pp.~101--103;
Stallings, Ch.~3 \S3.1, p.~108.}}

A useful image: a program is a recipe written down; a process is somebody
cooking it --- with the oven on, the ingredients half-mixed, and a finger on
line~7. Table~\ref{tab:prog-proc} makes the contrast precise.

\begin{table}[h]
\centering\small
\renewcommand{\arraystretch}{1.3}
\begin{tabular}{@{}>{\raggedright\arraybackslash}p{0.42\linewidth}>{\raggedright\arraybackslash}p{0.42\linewidth}@{}}
\toprule
\textbf{Program} & \textbf{Process} \\
\midrule
A file on disk (an executable) & A program in execution \\
\emph{Passive}: a list of instructions & \emph{Active}: has a program counter and live resources \\
One copy, however many runs & One per run, each independent \\
Exists until deleted & Created, runs, and terminates \\
\bottomrule
\end{tabular}
\caption{Program versus process (Silberschatz et al., Ch.~3 \S3.1, pp.~101--103).}
\label{tab:prog-proc}
\end{table}

A program becomes a process when its executable file is \key{loaded into
memory}. Two processes may run the same program; they are nevertheless two
separate execution sequences.

\begin{definitionbox}{Process}
A process is a program in execution: the unit of work in a modern time-sharing
system. It consists of the program code together with its current activity ---
the value of the program counter, the contents of the processor's registers,
and its memory.
\end{definitionbox}

Batch systems called the same entity a \key{job}, and the textbook uses
\emph{job} and \emph{process} almost interchangeably (Silberschatz, p.~101).
Stallings (\S3.1, p.~108) lists equivalent views: a program in execution, an
instance of a program running on a computer, and the entity that can be
assigned to and executed on a processor.

\begin{example}[\texttt{./counter} launched twice]
\begin{enumerate}
  \item Terminal A runs \texttt{./counter}. The OS creates \texttt{P1}: it
        loads the code, allocates a stack, and sets the program counter to the
        entry point.
  \item Terminal B runs \texttt{./counter}. The OS creates \texttt{P2} from
        the \emph{same} file, with a \emph{separate} stack and data.
  \item After a few seconds \texttt{P1} has printed 812 and \texttt{P2} has
        printed 305.
  \item The variable \texttt{count} therefore has two values at once --- one in
        each process's data.
  \item Both program counters may point into the same loop, at different
        instructions.
\end{enumerate}
\emph{Conclusion:} same text section, different activity; the OS tracks
\texttt{P1} and \texttt{P2} as two processes.
\end{example}

The counting criterion is not what is visible on screen. If you open three
windows of the same browser, the number of \emph{windows} does not decide how
many processes exist; the number of \emph{execution sequences the OS tracks}
does. One process can draw many windows, and one program can also start several
processes. The question to put to the OS is: how many PCBs exist for this
program?

\section{A Process in Memory}

\noindent\muted{\textit{Reading: Silberschatz et al., Ch.~3 \S3.1, pp.~101--102;
Stallings, Ch.~3 \S3.3, p.~126.}}

Before the kernel can record a process, the process needs a place to live. Its
address space has four regions, each with a job (Figure~\ref{fig:layout}).
The \key{text} section holds the program code. The \key{data} section holds
global variables. The \key{heap} holds memory allocated at run time and grows
upward. The \key{stack} holds function parameters, return addresses, and local
variables, and grows downward. Between heap and stack lies free space into
which both can grow.

\begin{figure}[h]
\centering
\begin{tikzpicture}[every node/.style={font=\small}]
  \draw[draw=primary-blue, fill=light-bg] (0,3.6) rectangle (3.2,4.4);
  \node at (1.6,4.0) {stack};
  \draw[draw=neutral, dashed] (0,2.4) rectangle (3.2,3.6);
  \draw[draw=primary-blue, fill=light-bg] (0,1.6) rectangle (3.2,2.4);
  \node at (1.6,2.0) {heap};
  \draw[draw=primary-blue, fill=light-bg] (0,0.8) rectangle (3.2,1.6);
  \node at (1.6,1.2) {data};
  \draw[draw=primary-blue, fill=light-bg] (0,0) rectangle (3.2,0.8);
  \node at (1.6,0.4) {text};
  \draw[observed-edge] (0.8,2.5) -- (0.8,3.0);
  \draw[observed-edge] (2.4,3.5) -- (2.4,3.0);
  \node[anchor=east, text=neutral] at (-0.25,4.4) {max};
  \node[anchor=east, text=neutral] at (-0.25,0) {0};
  \node[anchor=west, text=primary-gold, font=\footnotesize] at (3.55,4.0) {locals, returns};
  \node[anchor=west, text=primary-gold, font=\footnotesize] at (3.55,3.0) {free};
  \node[anchor=west, text=primary-gold, font=\footnotesize] at (3.55,2.0) {\texttt{malloc}};
  \node[anchor=west, text=primary-gold, font=\footnotesize] at (3.55,1.2) {globals};
  \node[anchor=west, text=primary-gold, font=\footnotesize] at (3.55,0.4) {code};
\end{tikzpicture}
\caption{The four regions of a process's address space, from address 0 (text)
up to the maximum address (stack); the heap grows up and the stack grows down
into the free band --- cf.\ Silberschatz et al., Ch.~3 \S3.1, pp.~101--102
\cite{SilberschatzOSC_operating_system_concepts}.}
\label{fig:layout}
\end{figure}

\begin{example}[Where does each variable live?]
Consider this fragment of \texttt{counter.c}:
\begin{verbatim}
int count = 0;  /* (a) */
int main(void) {
  int step = 1;  /* (b) */
  int *log =
    malloc(64);  /* (c) */
  while (1) {
    count += step; /* (d) */
    print(count);
  }
}
\end{verbatim}
(a) \texttt{count} is a global, so it lives in \key{data}. (b) \texttt{step} is
a local variable of \texttt{main}, so it lives on the \key{stack}. (c) The 64
bytes that \texttt{log} points to were allocated at run time, so they are on the
\key{heap}; the pointer variable \texttt{log} itself is a local and lives on the
\key{stack}. (d) The instructions of the loop are code, in \key{text}. For
\texttt{P1} and \texttt{P2} the text is the same, while data, heap, and stack
are separate: \texttt{count} is 812 in one and 305 in the other.
\end{example}

Stallings (\S3.3, p.~126) calls the whole collection a \key{process image}: the
user program (text), user data (data and heap), the stack, and the process's
\key{control block}, the kernel's record of it. The first three are what the
process uses; the last is what the \emph{kernel} uses to manage it. Put another
way, the memory regions describe what the process \emph{is doing}, while the
control block describes what the kernel \emph{knows about it}.

\section{Process States}

\noindent\muted{\textit{Reading: Silberschatz et al., Ch.~3 \S3.1, p.~103;
Stallings, Ch.~3 \S3.2, p.~111.}}

A process is not always running. Most of the time it is waiting --- for the CPU
or for the outside world --- and the kernel must know which. Silberschatz names
five states (Table~\ref{tab:states}).

\begin{table}[h]
\centering\small
\renewcommand{\arraystretch}{1.25}
\begin{tabular}{@{}>{\raggedright\arraybackslash}p{0.18\linewidth}>{\raggedright\arraybackslash}p{0.7\linewidth}@{}}
\toprule
\textbf{State} & \textbf{Meaning} \\
\midrule
new & The process is being created. \\
ready & The process is waiting to be assigned to a processor. \\
running & Instructions are being executed. \\
waiting & The process is waiting for some event (I/O completion, a signal). Stallings: \emph{blocked}. \\
terminated & The process has finished execution. Stallings: \emph{exit}. \\
\bottomrule
\end{tabular}
\caption{The five process states.}
\label{tab:states}
\end{table}

\begin{definitionbox}{Counting rule}
On one processor, at most one process is running at any instant; many may be
ready or waiting. With $k$ CPUs, at most $k$ processes run.
\end{definitionbox}

Figure~\ref{fig:states} shows the state diagram. It has exactly six
transitions, and each is caused by a specific event, listed in
Table~\ref{tab:trans}.

\begin{figure}[h]
\centering
\begin{tikzpicture}[every node/.style={font=\small},
    st/.style={dag-node, minimum width=2.3cm, minimum height=0.8cm}]
  \node[st] (new) at (-5.2,1.6) {new};
  \node[st] (ready) at (-2.0,0) {ready};
  \node[st] (run) at (2.0,0) {running};
  \node[st] (term) at (5.2,1.6) {terminated};
  \node[st] (wait) at (0,-2.2) {waiting};
  \draw[observed-edge] (new) -- node[above, sloped, font=\footnotesize] {admitted} (ready);
  \draw[observed-edge] (ready.north east) to[bend left=30]
    node[above, font=\footnotesize] {scheduler dispatch} (run.north west);
  \draw[observed-edge] (run.south west) to[bend left=30]
    node[below, font=\footnotesize] {interrupt} (ready.south east);
  \draw[observed-edge] (run.south) |- node[pos=0.25, right, font=\footnotesize] {I/O or event wait} (wait.east);
  \draw[observed-edge] (wait.west) -| node[pos=0.75, left, font=\footnotesize] {I/O or event completion} (ready.south);
  \draw[observed-edge] (run) -- node[above, sloped, font=\footnotesize] {exit} (term);
\end{tikzpicture}
\caption{The process-state diagram, after Silberschatz et al., Ch.~3 \S3.1,
p.~103 \cite{SilberschatzOSC_operating_system_concepts}.}
\label{fig:states}
\end{figure}

\begin{table}[h]
\centering\small
\renewcommand{\arraystretch}{1.25}
\begin{tabular}{@{}>{\raggedright\arraybackslash}p{0.28\linewidth}>{\raggedright\arraybackslash}p{0.62\linewidth}@{}}
\toprule
\textbf{Transition} & \textbf{Caused by} \\
\midrule
new $\to$ ready & the OS admits the process (resources allocated) \\
ready $\to$ running & the scheduler dispatches it to the CPU \\
running $\to$ ready & an interrupt, e.g.\ the timer: its time is up \\
running $\to$ waiting & it requests I/O or waits for an event (a blocking system call such as \texttt{read}) \\
waiting $\to$ ready & the I/O completes or the event occurs \\
running $\to$ terminated & it calls \texttt{exit} (or is killed) \\
\bottomrule
\end{tabular}
\caption{The cause of each transition.}
\label{tab:trans}
\end{table}

Two arrows are notable for being absent. There is no waiting $\to$ running
arrow: a waiting process must first become ready. And there is no ready $\to$
waiting arrow: a ready process is not executing, so it cannot request I/O. The
first absence deserves a reason. When a waiting process receives its I/O,
another process is probably running on the CPU right now; only the scheduler
decides who runs next, and the I/O completion only makes the process
\emph{eligible}. Keeping ``eligible'' and ``chosen'' separate is what lets the
OS apply a scheduling policy at all.

\begin{example}[\texttt{P1} reads a key]
\begin{enumerate}
  \item \texttt{P1} is created: \emph{new}.
  \item The OS admits it: \emph{ready}.
  \item The scheduler dispatches it: \emph{running}; it increments \texttt{count}.
  \item The timer interrupt fires: back to \emph{ready}; \texttt{P2} runs.
  \item \texttt{P1} is dispatched again: \emph{running}. It calls \texttt{read}
        on the keyboard: \emph{waiting}.
  \item The user presses a key; the interrupt completes the I/O: \emph{ready}.
  \item \texttt{P1} is dispatched: \emph{running}. It calls \texttt{exit}:
        \emph{terminated}.
\end{enumerate}
The trace visits seven states and takes six transitions. The blocking
\texttt{read} is the system call of Lecture~3 seen from the process's side.
\end{example}

\begin{example}[Counting states on one CPU]
Five processes exist. \texttt{P1} is running; \texttt{P2} and \texttt{P3} wait
for the disk; \texttt{P4} and \texttt{P5} are ready. \texttt{P1}'s time slice
expires. The interrupt moves \texttt{P1} to ready, and the scheduler picks
\texttt{P4}, which becomes running. Afterwards: running 1 (\texttt{P4}); ready
2 (\texttt{P5}, \texttt{P1}); waiting 2 (\texttt{P2}, \texttt{P3}). The total
is still five: transitions move processes, they do not create them.
\end{example}

Stallings (\S3.2, p.~111) reaches the same five states from a simpler start.
His \key{two-state} model has only \emph{running} and \emph{not running}. It is
too coarse: a not-running process might be ready to go or blocked on I/O, and
dispatching a blocked process wastes the CPU. Splitting \emph{not running} into
\emph{ready} and \emph{blocked}, and adding \emph{new} and \emph{exit}, gives
the \key{five-state} model. Stallings later adds \emph{suspended} states for
swapped-out processes (a seven-state model); we return to these with the
medium-term scheduler.

\section{The Process Control Block}

\noindent\muted{\textit{Reading: Silberschatz et al., Ch.~3 \S3.1, pp.~103--104;
Stallings, Ch.~3 \S3.3, p.~126; Bovet \& Cesati, Ch.~3, pp.~79--130.}}

When \texttt{P1} stops running, everything it was doing must be written down
somewhere --- or it can never be resumed. That somewhere is the PCB.

\begin{definitionbox}{Process control block}
The process control block (PCB), also called a task control block, is the
kernel data structure that represents one process. It is the repository for any
information that may vary from process to process.
\end{definitionbox}

There is one PCB per process: \texttt{P1} and \texttt{P2} have two PCBs even
though they share one program file. The PCB lives in \emph{kernel} memory, and
the process cannot read or write it directly. Creating a process means creating
a PCB; terminating it means eventually freeing the PCB. Figure~\ref{fig:pcb}
shows its principal fields.

\begin{figure}[h]
\centering
\begin{tikzpicture}[every node/.style={font=\small},
    fld/.style={draw=primary-blue, fill=light-bg, minimum width=4.6cm,
      minimum height=0.6cm, inner sep=1pt}]
  \node[text=primary-gold, font=\small\bfseries] at (0,0.6) {PCB of \texttt{P1}};
  \node[fld] at (0,0) {process state};
  \node[fld] at (0,-0.6) {process number (PID)};
  \node[fld] at (0,-1.2) {program counter};
  \node[fld] at (0,-1.8) {registers};
  \node[fld] at (0,-2.4) {memory limits};
  \node[fld] at (0,-3.0) {list of open files};
  \node[fld] at (0,-3.6) {$\cdots$};
\end{tikzpicture}
\caption{The fields of a PCB, after Silberschatz et al., Ch.~3 \S3.1, p.~103
\cite{SilberschatzOSC_operating_system_concepts}.}
\label{fig:pcb}
\end{figure}

The \emph{state} is one of new, ready, running, waiting, terminated. The
\emph{PID} is the unique process number. The \emph{program counter} holds the
address of the next instruction to execute. The \emph{registers} are the
accumulators, index registers, stack pointer, general-purpose registers, and
condition codes. The \emph{memory limits} say where the process's regions are,
and the \emph{open files} say what it is using. Silberschatz (pp.~103--104)
groups the information into seven categories (Table~\ref{tab:pcbcat}).

\begin{table}[h]
\centering\small
\renewcommand{\arraystretch}{1.25}
\begin{tabular}{@{}>{\raggedright\arraybackslash}p{0.3\linewidth}>{\raggedright\arraybackslash}p{0.6\linewidth}@{}}
\toprule
\textbf{Category} & \textbf{Examples} \\
\midrule
Process state & new, ready, running, waiting, halted \\
Program counter & address of the next instruction \\
CPU registers & accumulators, index registers, stack pointers, general-purpose registers, condition codes \\
CPU-scheduling information & priority, pointers to scheduling queues \\
Memory-management information & base and limit registers, page tables or segment tables \\
Accounting information & CPU time used, time limits, job or process numbers \\
I/O status information & I/O devices allocated, list of open files \\
\bottomrule
\end{tabular}
\caption{The seven categories of PCB information (Silberschatz et al., Ch.~3
\S3.1, pp.~103--104).}
\label{tab:pcbcat}
\end{table}

Why save the program counter and registers in particular? When an interrupt
takes the CPU away from \texttt{P1}, the CPU's registers still hold
\texttt{P1}'s values, and \texttt{P2} is about to overwrite them. So the kernel
first copies \texttt{P1}'s PC and registers into \texttt{P1}'s PCB. Later it
copies them back, and \texttt{P1} continues as if nothing had happened. State
information must be saved when an interrupt occurs, to allow the process to be
continued correctly afterward; the PCB is what makes interruption safe. Saving
one process's state and loading another's is the \key{context switch}, the
topic of the next lecture.

\begin{example}[Two PCBs, one program]
At one instant the two \texttt{./counter} processes look like this:

\begin{center}\small
\renewcommand{\arraystretch}{1.2}
\begin{tabular}{@{}lll@{}}
\toprule
\textbf{Field} & \textbf{PCB of \texttt{P1}} & \textbf{PCB of \texttt{P2}} \\
\midrule
state & waiting (keyboard) & running \\
PID & 4120 & 4127 \\
program counter & inside \texttt{read} wrapper & inside the loop \\
\texttt{count} (in data) & 812 & 305 \\
open files & terminal A & terminal B \\
text section & \multicolumn{2}{c}{shared: same \texttt{./counter} code} \\
\bottomrule
\end{tabular}
\end{center}
Every row that differs is per-process information --- exactly what the PCB
exists to hold. The PIDs are illustrative; the OS picks them.
\end{example}

Stallings (\S3.3, p.~126) groups the same content by purpose into three parts.
\emph{Process identification}: PID, parent's PID, user ID. \emph{Processor state
information}: user-visible registers, control and status registers (PC, flags),
and stack pointers. \emph{Process control information}: scheduling and state,
links to other processes, privileges, and memory and resource ownership. The
processor-state part is what a context switch saves and restores.

In Linux the PCB is the \key{process descriptor}, a C structure of type
\texttt{task\_struct} (Bovet \& Cesati, Ch.~3, pp.~79--130
\cite{BovetCesati_understanding_linux_kernel}). It holds the state, the PID,
scheduling information, a pointer to the memory descriptor, the open-file table,
parent and child links, and much more. Its \texttt{state} field records values
such as \texttt{TASK\_RUNNING}, \texttt{TASK\_INTERRUPTIBLE}, and
\texttt{TASK\_UNINTERRUPTIBLE}, and the kernel links all descriptors into a list
so it can walk every process. Note the naming: Linux uses
\texttt{TASK\_RUNNING} for both ready and running, so the textbook's two states
map to one flag plus the question ``is it on a CPU right now?''

Why can the PCB not be stored in the process's own stack or heap, where it
would be cheaper to reach? Because the process could then overwrite its own
state, PID, or memory limits and break protection; and because the kernel must
read the PCB even when the process is not running and its memory may be swapped
out. So PCBs live in kernel memory, reachable only in kernel mode.

\section{Threads: A Preview}

\noindent\muted{\textit{Reading: Silberschatz et al., Ch.~3 \S3.1, pp.~104--105;
threads in full in Ch.~4.}}

A process is its address space plus its PCB. One more refinement completes the
picture. The model so far assumes that a process performs a \emph{single thread
of execution} --- one program counter. Modern systems let a process have
\key{multiple threads}, each with its own PC and registers, sharing the same
address space; the PCB is then extended with per-thread information. A word
processor with one thread cannot type characters and run the spell checker at
the same time; with two threads it can.

\section{Misconceptions and Summary}

Four errors recur. ``A process is the program file'' --- no: a process is a
program \emph{in execution}. ``Two runs of one program share one PCB'' --- no:
there is one PCB per process, even for the same program. ``I/O completion puts a
process back on the CPU'' --- no: it goes to \emph{ready}, and the scheduler
decides. ``Many processes can be running on one CPU'' --- no: at most one at a
time per CPU; the others are ready or waiting.

In summary: a \key{process} is a program in execution; the program is passive
and the process is active. Its address space has \key{text}, \key{data},
\key{heap}, and \key{stack}, and two processes running one program share text
only. It moves among five \key{states} --- new, ready, running, waiting,
terminated --- through six legal transitions, each caused by an event. The
\key{PCB} records everything per-process: state, PC, registers, scheduling,
memory, accounting, and I/O status. Saving the PC and registers into the PCB is
what makes interruption and resumption possible.

\section{Looking Ahead and Further Reading}

The next lecture examines the queues that hold PCBs, the long-, short-, and
medium-term schedulers that move them, and the context switch that saves one
PCB and loads another. Its reading is Silberschatz et al., Ch.~3 \S3.2
\cite{SilberschatzOSC_operating_system_concepts}; Stallings, Ch.~3 \S3.2 and
\S3.4 \cite{StallingsOS_internals_design_principles}; and Bovet \& Cesati,
Ch.~3 \cite{BovetCesati_understanding_linux_kernel}. For this chapter the
page-cited sources are:
\begin{itemize}
  \item Silberschatz, Galvin, and Gagne, \emph{Operating System Concepts}, 8th
        ed. --- Ch.~3 \S3.1 (the process, p.~101; process state, p.~103;
        process control block, p.~103; threads, pp.~104--105).
  \item Stallings, \emph{Operating Systems: Internals and Design Principles},
        6th ed. --- Ch.~3 \S3.1 (p.~108), \S3.2 (p.~111), \S3.3 (p.~126).
  \item Bovet and Cesati, \emph{Understanding the Linux Kernel}, 3rd ed. ---
        Ch.~3 (pp.~79--130), for the concrete \texttt{task\_struct}.
\end{itemize}
All are page-verified through each book's \texttt{index.md} in the course supporting-books directory:
\begin{center}\texttt{master\_supporting\_docs/CS502/supporting\_books/}\end{center}

\section*{Self-Check}
\addcontentsline{toc}{section}{Self-Check}

\subsection*{Exercises}
\begin{enumerate}
  \item State, for each of the following, which region of the address space
        holds it: a \texttt{static} global array, the return address of a
        function call, a buffer returned by \texttt{malloc}, the machine code of
        \texttt{main}.
  \item A process goes: created, admitted, dispatched, timer interrupt,
        dispatched, requests disk I/O, I/O completes, dispatched, exits. List
        the states it visits and count the transitions.
  \item On a dual-core machine, eight processes exist: three are blocked on
        I/O. What is the maximum number running, and the minimum number ready?
  \item Explain why a waiting process cannot move directly to running.
  \item Which PCB category holds each of: the page-table pointer, the priority,
        the CPU time used, the list of open files?
\end{enumerate}

\subsection*{Solutions}
\begingroup\small
\begin{enumerate}
  \item Global array: data. Return address: stack. \texttt{malloc} buffer:
        heap. Machine code of \texttt{main}: text.
  \item States: new, ready, running, ready, running, waiting, ready, running,
        terminated --- nine states visited (counting repeats), and eight
        transitions: admitted, dispatch, interrupt, dispatch, I/O wait, I/O
        completion, dispatch, exit.
  \item At most $k=2$ are running. The five non-blocked processes are either
        running or ready, so with two running at least $5-2=3$ are ready.
  \item I/O completion only makes the process eligible; another process is
        probably running, and only the scheduler chooses which ready process
        runs. Separating eligibility from selection is what allows a
        scheduling policy.
  \item Page-table pointer: memory-management information. Priority:
        CPU-scheduling information. CPU time used: accounting information. Open
        files: I/O status information.
\end{enumerate}
\endgroup

\bibliographystyle{plain}
\bibliography{../../Bibliography_base}

\end{document}
